\documentclass[a4,11pt]{aleph-notas}

% -- Paquetes adicionales
\usepackage{enumitem}
\usepackage{aleph-comandos}

% -- Datos
\institucion{Facultad de Ciencias Exactas, Naturales y Ambientales}
\carrera{Catálogo STEM}
\asignatura{Matemática III}
\tema{Resumen 08: Función implícita}
\autor[A. Merino]{Andrés Merino}
\fecha{Periodo 2026-1}

\logouno[0.16\textwidth]{Logos/logoPUCE_04_ac}
\definecolor{colortext}{HTML}{1957d1}
\definecolor{colordef}{HTML}{1957d1}
\fuente{montserrat}


\begin{document}

\encabezado

En este resumen consideraremos $\Omega\subset\R^n$ un conjunto abierto.

%%%%%%%%%%%%%%%%%%%%%%%%%%%%%%%%%%%%%%
\section{Función implícita}
%%%%%%%%%%%%%%%%%%%%%%%%%%%%%%%%%%%%%%

\begin{defi}
    Sean $\func{f}{\Omega}{\R}$ y $\func{\hat{f}}{\R^{n+1}}{\R}$ campos escalares. Se dice que $f$ está definida de manera implícita con respecto a $\hat{f}$ si para todo $(x_1,\ldots,x_n)\in \Omega$ se tiene que
    \[
        \hat{f}(x_1,\ldots,x_n,f(x_1,\ldots,x_n))=0.
    \]
\end{defi}

\begin{teo}
    Sean $\func{f}{\Omega}{\R}$ y $\func{\hat{f}}{\R^{n+1}}{\R}$ campos escalares tales que $f$ está definida de manera implícita con respecto a $\hat{f}$. Si $\hat{f}$ es diferenciable, entonces para cada $k=1,\ldots, n$ se tiene que
    \[
        D_k f(x_1,\ldots,x_n) =
        -\frac{D_k \hat{f}(x_1,\ldots,x_n,f(x_1,\ldots,x_n))}{D_{n+1} \hat{f}(x_1,\ldots,x_n,f(x_1,\ldots,x_n))}
    \]
    para todo $(x_1,\ldots,x_n)\in \Omega$ tal que $D_{n+1} \hat{f}(x_1,\ldots,x_n,f(x_1,\ldots,x_n))\neq 0$.
\end{teo}

\begin{advertencia}
    Bajo el teorema anterior,
    \[
        \parcial{f}{x_k}(x_1,\ldots,x_n)=
        -\frac{\parcial{\hat{f}}{x_k}(x_1,\ldots,x_n,f(x_1,\ldots,x_n))}
              {\parcial{\hat{f}}{x_{n+1}}(x_1,\ldots,x_n,f(x_1,\ldots,x_n))}.
    \]
\end{advertencia}


\end{document}
